The reference workflow is intentionally centralized so that citation state has a single owner. The shared bibliography file, \texttt{references/references.bib}, is maintained by \texttt{references\_agent} rather than by individual section agents. That division of responsibility prevents concurrent edits from fragmenting the bibliography, keeps citation keys canonical, and gives the system one place to register, validate, and update reference records.

Section agents therefore treat the bibliography as read-only support data. When drafting or revising a section, a section agent first checks the registered reference context and verifies that every citation key it intends to use already exists in \texttt{references/references.bib}. If a key is registered, the section can cite it directly with the house citation command or \texttt{\textbackslash cite\{key\}} as appropriate. If a key is not registered, the section agent does not invent a citation, guess at a source, or silently omit the support. Instead, it routes the gap into the reference workflow by requesting research or by sending a structured handoff to \texttt{references\_agent}.

That handoff is the mechanism that keeps unsupported claims visible. A missing citation, missing definition, or otherwise ungrounded statement becomes an explicit task rather than an implicit assumption. The section agent can queue a \texttt{do\_research\_on\_the\_web} task when the needed source is not yet known, then forward candidate bibliographic records to \texttt{references\_agent} for registration. Once \texttt{references\_agent} confirms the new keys, the section agent verifies that those keys are present in \texttt{references/references.bib} and only then revises the prose to use them.

This ownership model has two practical benefits. First, it keeps citation support auditable: every key used in a section can be traced back to a registered bibliography entry. Second, it keeps section drafting focused on prose rather than bibliography maintenance. The section agent is responsible for recognizing where support is needed and for integrating approved keys into the text, while \texttt{references\_agent} is responsible for the canonical reference store itself. In the Codynamic Book Machine, that separation is what turns citation management into a durable workflow instead of an ad hoc editing habit.

The same pattern also clarifies how reference support fits into the broader chapter. Section agents do not merely ask whether a claim sounds plausible; they check whether the manuscript already has the support needed to state it responsibly. If the support exists, they cite it. If it does not, they route the gap to the specialist workflow that can resolve it. That discipline keeps the manuscript honest about what is grounded in registered evidence and what still requires follow-up.

In the surrounding chapter, this section establishes the reference side of specialist handoffs. The next section extends the same pattern to diagrams: section agents may identify a visual need, but they must route the request through the diagram workflow and wait for a callback before inserting any asset path into the section payload.
